% Optional continuation of full_gaussian_rank.tex; same theorem environments.
\section{A universal splitting at every cyclotomic level}
\label{fg:sec:universal}

The boundary lift used above does not depend on the number four. It
separates the same-point coordinates from the remaining color problem
at every cyclotomic level.

Let $N\ge3$, $\omega=e^{2\pi i/N}$, and define the formal imaginary
coordinates for $D_{a,w-a}(\omega^r,\omega^s)$ exactly as before, now
with colors modulo $N$. Impose conjugation, discard its fixed
coordinates, and remove one divergent coordinate $X_1^{0,r}$ for
each conjugate pair of nonreal roots. Let $A_{N,w}$ be the matrix of
finite single-product shuffle and stuffle rows, with only their
difference retained when exactly one factor diverges. Put
\[
 R_N=\{1,\ldots,\lfloor(N-1)/2\rfloor\},\qquad k_N=|R_N|.
\]
The matrix has
\[
 d_{N,w}=\frac{N^2-\gcd(N,2)^2}{2}(w-1)-k_N
\]
columns: conjugation has $\gcd(N,2)^2$ fixed color pairs, and the
last term counts the removed divergent coordinates. Let
$A_{N,w}^{\neg g}$ delete all $k_N(w-1)$ same-point columns
$X_a^{r,0}$ with $r\in R_N$. These are formal matrices with the same
restricted meaning as $A_w$; distribution and higher-depth rows
have not been added.

For a homogeneous degree-$n$ polynomial, where $n=w-2$, define
\[
 \mathcal H_n^-=
 \{G\in\mathbb Q[X,Y]_n:G(X,Y)=-G(X,X-Y)\}.
\]

\begin{theorem}[Universal same-point splitting]\label{fg:thm:universal}
For odd $w=2m+1$, there is a canonical split exact sequence
\[
 0\longrightarrow\ker A_{N,w}^{\neg g}
 \longrightarrow\ker A_{N,w}
 \longrightarrow\bigoplus_{r\in R_N}\mathcal H_{2m-1}^-
 \longrightarrow0.
\]
The last map records the same-point polynomials $G_r$.
In particular,
\begin{equation}\label{fg:eq:universal-nullity}
 \dim\ker A_{N,w}-\dim\ker A_{N,w}^{\neg g}=k_Nm.
\end{equation}
At even $w$, every vector in $\ker A_{N,w}$ has all same-point
coordinates zero, so the two kernels are naturally identical.
Consequently
\begin{equation}\label{fg:eq:universal-rank-difference}
 \operatorname{rank}A_{N,w}-\operatorname{rank}A_{N,w}^{\neg g}
 =\begin{cases}
    k_N(w-1),&w\text{ even},\\
    k_N(w-1)/2,&w\text{ odd}.
  \end{cases}
\end{equation}
\end{theorem}
\begin{proof}
Fix $r\in R_N$ and abbreviate the polynomials for colors
$(r,0),(0,r),(r,-r)$ by $G_r,U_r,M_r$. The product-color classes
$(r,r),(0,r),(r,-r)$ imply exactly the same necessary conditions
as in the Gaussian proof:
\begin{align*}
 G_r(X,Y)&=-G_r(X,X-Y),\\
 U_r(X,Y)&=-G_r(Y,X)+c_rY^n,\qquad c_r=[X^n]G_r,\\
 M_r(X,Y)&=G_r(X-Y,X),\qquad M_r=M_r^\sigma.
\end{align*}
The boundary equation uses the single regularized difference,
so its derivation is unchanged. Homogeneity gives
$M_r^\sigma=(-1)^{n+1}M_r$. Thus even $n$ forces $G_r=0$.
For odd $n$, the symmetry of $M_r$ is automatic.

For odd $n$, define a section by assigning the displayed values to
these three color pairs for each $r$, their negatives to conjugate
pairs, and zero to every other color pair. This assignment satisfies
all rows. Here is a support check that prevents an implicit assumption
about the other colors. A shuffle can meet a nonzero $G_r$ coordinate
only when its two product colors are $(r,r)$, and then it is exactly
the antireflection equation. It can meet $U_r$ or $M_r$ only for
product colors $(0,r)$ or $(r,0)$, and then it is exactly the boundary
equation. The stuffle rows meeting these coordinates are the boundary
class and $(r,-r)$, which is the symmetry of $M_r$. Transposition
and conjugation cover the reversed and conjugate classes. Thus no
unexamined row couples the proposed section to other colors.
The supports for distinct $r\in R_N$ are disjoint.

Subtracting this section from any kernel vector leaves all $G_r=0$;
the remainder is naturally a vector in $\ker A_{N,w}^{\neg g}$.
This proves the splitting. Since
$\dim\mathcal H_{2m-1}^-=m$, it gives
\eqref{fg:eq:universal-nullity}. Subtracting kernel dimensions from
column counts proves \eqref{fg:eq:universal-rank-difference}.
\end{proof}

At odd weight, an immediate consequence is that every coefficient
relation supported only on these same-point columns is spanned by
the independent same-point shuffle rows, separately for each root
pair. Those rows already supply $k_N(w-1)/2$ dimensions, equal to
\eqref{fg:eq:universal-rank-difference}. Thus the Gaussian-only
completeness result above is one instance of a level-independent
fact about this particular single-product coefficient system.

\subsection{The complete level-three ranks}

The earlier incoming report \cite{PriorResearch} already identifies the residual
level-three space after $G=0$. In that case the boundary formulas
force $U=M=0$, so its regularized presentation and the present
removed-coordinate presentation agree. Its remaining polynomial
satisfies
\[
 A(Y,X)=-A(X,Y),\qquad A(X,X-Y)=A(X,Y),
\]
and is parameterized by
\[
 A=\Delta\,\mathbb Q[Q,\Delta^2],\qquad
 \Delta=XY(X-Y),\quad Q=X^2-XY+Y^2.
\]
At odd weight its homogeneous degree-$w-2$ component has dimension
$\lfloor(w+1)/6\rfloor$, and at even weight it is zero.
For completeness, antisymmetry and reflection force the three
linear factors of $\Delta$; after their removal the quotient is
invariant under the group generated by the two transformations.
That group acts as permutations and global negation on
$(X,-Y,Y-X)$, whose sum is zero, giving the invariant ring
$\mathbb Q[Q,\Delta^2]$. Counting monomials of degrees
$3+2a+6b=w-2$ gives the stated dimension.
This residual description is credited to the preceding report; the
new assertion is its extension to the free same-point family under
the manuscript's boundary convention.

\begin{corollary}[Level-three extension]\label{fg:cor:level3}
The analogous full level-three matrix has $4w-5$ columns and satisfies
\[
 \operatorname{rank}A_{3,w}
 =\begin{cases}
    4w-5,&w\text{ even},\\
    (7w-9)/2-\lfloor(w+1)/6\rfloor,&w\text{ odd},
  \end{cases}
\]
while the matrix with all same-point columns deleted satisfies
\[
 \operatorname{rank}A_{3,w}^{\neg g}
 =\begin{cases}
    3w-4,&w\text{ even},\\
    3w-4-\lfloor(w+1)/6\rfloor,&w\text{ odd}.
  \end{cases}
\]
\end{corollary}
\begin{proof}
At odd weight the full kernel has dimension
$(w-1)/2+\lfloor(w+1)/6\rfloor$ by
Theorem~\ref{fg:thm:universal}; its restricted kernel has only the
second summand. Both vanish at even weight. Subtraction from the
respective column counts proves the formulas.
\end{proof}

The supplementary program \texttt{universal\_gaussian\_lift.py}
verifies the lift directly against every raw row at selected weights
and levels $3,4,5,6,8,10,12$. Independent exact ranks in weights
$2,3,4$ at levels $3$ through $6$ check the rank differences. This
establishes neither a formula for the remaining kernel at arbitrary
level nor numerical period independence; those remain separate
questions.
